\chapter{Future Work}
\label{chap:future-work}

\section{External and longitudinal validation}

The immediate priority is an external validation study using independently acquired MAIA and OCT data. It should include more than one clinical site where possible and should record scanner type, acquisition protocol and image quality so that any loss of performance can be interpreted. A prospective design would also allow the eligibility criteria, acceptable agreement range, calibration assessment and handling of missing or low-quality examinations to be specified before model evaluation. Such a study should retain participant-level splits and report performance separately across clinically relevant disease stages rather than relying on a pooled point-level result.

Longitudinal data are particularly important. The RUSH2A natural-history studies have shown that both OCT measures and microperimetry change over time, while functional transition-point sensitivity may be responsive to clinically meaningful decline \citep{Vincent2025,Maguire2024}. Future work should therefore test whether image-derived estimates track within-person change across repeated visits, not only cross-sectional differences between retinal points. This will require a model that accounts for repeated measures and for the fact that deterioration may be slow and heterogeneous. Pointwise sensitivity may remain useful, but mean sensitivity and sensitivity at carefully defined transition regions should be evaluated alongside it when the intended use is trial monitoring.

\section{Outcome, uncertainty and spatial representation}

Repeated microperimetry at the same visits would make it possible to model measurement reliability directly. In particular, a future analysis could distinguish observations at the device floor, quantify test--retest variation in the study population and report prediction intervals that combine model uncertainty with outcome variability. This would be more useful than a single MAE when deciding whether an estimate is sufficiently reliable for a particular research task \citep{Pfau2021Microperimetry,Karuntu2025}.

The spatial ablation in this thesis also suggests a focused methodological question. Multiple within-footprint A-scan positions should not simply be averaged because pooling smoothed the local representation and was less effective for the pointwise target in this cohort. With a larger dataset, future work could compare pre-specified alternatives: one centred anchor, learned attention across a small set of anchors, and representations that preserve the distribution of local features without collapsing it into a mean. These comparisons should be conducted within the same participant-grouped validation design. The three-B-scan summary should likewise be tested against explicit models of scan-to-scan variability, while maintaining the link between each OCT input and the MAIA stimulus location.

\section{Boundary data and representation development}

The B0--B4 source should be linked to an authoritative description of the boundary definitions before any named-layer interpretation is attempted. Once this provenance has been verified, a larger boundary-available cohort could compare engineered interval features, segmentation-derived features and boundary-aligned image representations on the same cases. Disease-specific segmentation is likely to be important: an ellipsoid-zone model trained in another disease transferred poorly to \textit{USH2A}-related degeneration until it was retrained on relevant data \citep{Loo2022}. The same principle applies to foundation-model features. RETFound provides a reasonable starting representation, but fine-tuning or further self-supervised adaptation should only be considered with enough independent data to evaluate whether it improves on the current constrained model \citep{Zhou2023RETFound,Gholami2024}.

\section{Path towards clinical use}

Any clinical or trial-facing tool would need to be evaluated as a decision-support intervention rather than only as a prediction model. Its intended user, input requirements, output, uncertainty display and failure cases should be defined in advance, together with evidence that using the output improves a real decision. These are consistent with the reporting priorities set out in CONSORT-AI, including transparent descriptions of inputs, human--AI interaction and error analysis \citep{Liu2020}. Until that evidence is available, the most appropriate role for OCT-based sensitivity estimates is as a research adjunct that complements, rather than replaces, direct functional testing.
